\section{A. The tower and the manuscript's changing definitions}
\subsection{A.1. Surfaces, punctures, and the two finite characters}

\statement{A.1}{Proved}{The permutation and orientation conventions}
Let $X_g$ be a closed connected orientable topological surface and $S_n\subset X_g$ a specified $n$-element set. Give its homeomorphism group the compact-open topology on a homeomorphism and its inverse, and put $U_{g,n}=X_g\setminus S_n$. Then
\[
1\longrightarrow A^!_{g,n}\longrightarrow A_{g,n}
 \longrightarrow\mathfrak S_n\longrightarrow1,
\qquad
1\longrightarrow\pi_0(A^!_{g,n})\longrightarrow M^{\pm}_{g,n}
 \longrightarrow\mathfrak S_n\longrightarrow1.
\]
The orientation character $\epsilon:M^{\pm}_{g,n}\to\{\pm1\}$ is surjective and its kernel is $M^+_{g,n}$. The intersection of the permutation and orientation kernels is $P_{g,n}$. In particular the manuscript's unadorned group is neither pure nor orientation-preserving by definition.
\Man{19}{102r--v}{204--205}; the positive superscript is specified in \S19, folio 113, scan 226.
\begin{proof}
Restriction to the finite set gives the first homomorphism. Every permutation is realized by an orientation-preserving homeomorphism: one realizes transpositions by interchanging two small discs along an embedded arc avoiding the other punctures, and composes. Its kernel is precisely the pointwise stabilizer. A path of homeomorphisms induces a constant permutation, so the identity components of the two groups coincide. Taking component groups therefore gives the second sequence, including surjectivity. Degree on the oriented fundamental class gives $\epsilon$. A reflection model for a genus-$g$ surface, followed by a point-moving homeomorphism, realizes degree $-1$ preserving the prescribed punctures individually. Thus both signs occur even in the pure group. No splitting of the permutation sequence is asserted.

Every homeomorphism of $U_{g,n}$ extends uniquely to its finite end compactification $X_g$: the ends are exactly the deleted punctures, and a homeomorphism permutes ends. This proves the abstract isomorphism in manuscript (5). The extension and restriction respect the indicated homeomorphism topologies: a compact core together with one smaller annular neighbourhood of each end controls the image and inverse image of the complementary discs. This provides uniform control on the compactification from compact-open control on the punctured surface.
\end{proof}

\statement{A.2}{Known}{Surface groups with peripheral structure}
Choose an orientation, a base point $b\in U_{g,n}$, and paths to small positively oriented puncture circles. Then
\[
\Pi_{g,n}=\left\langle a_1,b_1,\ldots,a_g,b_g,c_1,\ldots,c_n\ \middle|\
 \prod_{i=1}^g[a_i,b_i]\prod_{j=1}^nc_j=1\right\rangle.
\]
For $n>0$ this is free of rank $2g+n-1$. For hyperbolic type it is centre-free and residually finite. Its peripheral datum consists of the conjugacy classes of the cyclic subgroups $\langle c_j\rangle$, with the coherent orientation datum; for $n=0$ we retain the orientation line in $H_2(\Pi_{g,0},\mathbb Z)$. The Dehn--Nielsen--Baer identification is
\[
 M^{\pm}_{g,n}\xrightarrow{\sim}\Out_{\mathrm{lac}}(\Pi_{g,n}).
\]
The target preserves the peripheral classes as a set; restricting to fixed classes and multiplier $+1$ gives $P_{g,n}$.
\Man{19}{111--113}{221--226}; see \citet[Theorems 8.1 and 8.8]{AF_FarbMargalit2012}.
\begin{proof}[Verification and scope]
The presentation follows by cutting into a polygon and applying van Kampen. Eliminating $c_n$ proves the free-rank assertion. For a closed hyperbolic surface the universal covering is the disc and the deck group is torsion-free cocompact Fuchsian; its centre is trivial. Free groups of rank at least two have trivial centre by reduced words. Residual finiteness follows from the faithful finitely generated characteristic-zero linear realizations and reduction modulo finite residue fields. The identification with peripheral-preserving outer automorphisms is exactly the punctured or closed Dehn--Nielsen--Baer theorem cited above; this is an imported theorem, not an assertion that an arbitrary group automorphism preserves omitted punctures. Changing $b$ or paths conjugates the resulting fundamental-group maps. Completing yields $\Pi_{g,n}\hookrightarrow\pi_{g,n}$; it does not yet prove faithfulness of $\widehat{M^+_{g,n}}\to\Out(\pi_{g,n})$.
\end{proof}

\statement{A.3}{Known}{Filling one point and the precise Birman group}
For a hyperbolic target $U_{g,n}$ and one extra point $b$, let $M^{\pm}_{g,n+1,*}$ be the subgroup fixing that extra point, while allowing the old punctures to be permuted. There are exact sequences of discrete groups
\[
1\longrightarrow\Pi_{g,n}\xrightarrow{\mathrm{push}}
 M^{\pm}_{g,n+1,*}\xrightarrow{\mathrm{forget}}M^{\pm}_{g,n}\longrightarrow1,
\qquad
1\longrightarrow\Pi_{g,n}\longrightarrow P_{g,n+1}\longrightarrow P_{g,n}\longrightarrow1.
\]
There is no forgetting homomorphism on the full group that permits the extra point itself to be permuted. On fundamental groups the inclusion of opens instead gives the surjection
$\Pi_{g,n+1}\to\Pi_{g,n}$ killing the normal closure of the last peripheral loop. These are different maps between different groups.
\Man{19}{103--111}{206--221}; \citet[Theorem 4.6]{AF_FarbMargalit2012}.
\begin{proof}[Source and comparison maps]
Evaluation at $b$ is a locally trivial fibration of homeomorphism groups, by isotopy extension. Its homotopy sequence, and contractibility of the identity component in hyperbolic type, give the first exact sequence. The connecting map pushes $b$ along a based loop. Its conjugation action on $\Pi_{g,n}$ is the point-pushing action and its injectivity uses the centre-free hyperbolic hypothesis. Restricting to homeomorphisms fixing the old punctures gives the second sequence. Surjectivity follows by isotoping the image of $b$ back to $b$ through an arc missing the punctures. For the surface-group map, the filled disc attaches along the last boundary circle, so van Kampen gives exactly the indicated normal quotient. Section C distinguishes their profinite analogues and their kernels.
\end{proof}

\statement{A.4}{Proved}{A covering gives a correspondence with specified lifts}
Let $p:U'\to U$ be a connected finite unramified covering of orientable finite-type surfaces, with $U$ hyperbolic. A choice of points above $b$ identifies $\pi_1(U')$ with a finite-index subgroup $H<\Pi=\pi_1(U)$. Put $L_H=\Stab_{M^{\pm}(U)}([H])$, where $[H]$ denotes conjugacy in $\Pi$. The group of isotopy classes of pairs of commuting lifts $(f',f)$ fits into
\[
1\longrightarrow N_\Pi(H)/H\longrightarrow\mathrm{Lift}(p)
 \longrightarrow L_H\longrightarrow1,
\qquad \mathrm{Lift}(p)\longrightarrow M^{\pm}(U').
\]
Thus the general covering operation in the tower is a correspondence, not a canonical homomorphism from all of $M^{\pm}(U)$ to $M^{\pm}(U')$. It preserves the sign of orientation and the induced peripheral subgroups. The degree $d=[\Pi:H]$ satisfies
$2g'-2+n'=d(2g-2+n)$.
\Man{19 and 28}{102--111 and 177--180}{204--221 and 350--356}.
\begin{proof}
A base homeomorphism lifts if and only if its outer action preserves $[H]$, by the covering-space lifting criterion. Any two lifts differ by a deck transformation, and the deck group is $N_\Pi(H)/H$. Lifting an isotopy gives the asserted exact sequence on isotopy classes. A deck transformation isotopic to the identity on a hyperbolic surface is trivial: it has finite order, whereas a nontrivial periodic orientation-preserving mapping class cannot be represented by an identity-isotopic map; equivalently use its faithful action on the hyperbolic universal cover and the centre-free surface group. This also verifies that the kernel has not been accidentally collapsed. Forgetting the lower homeomorphism gives the displayed map to the upstairs mapping class group. The restriction of a representative automorphism to $H$ changes, when the lift changes, by the deck action. Even for a characteristic subgroup this is not generally an inner automorphism of $H$. The correspondence records that ambiguity. Peripheral annuli lift to peripheral annuli, and Euler characteristic multiplies under a finite covering; both assertions follow by lifting a finite cell decomposition. The same subgroup inclusion completes to an open embedding $\widehat H\hookrightarrow\widehat\Pi$, since a finite-index subgroup inherits the full profinite topology.
\end{proof}

\subsection{A.2. Teichm\"uller space and the modular topos}

\statement{A.5}{Known}{The modular object retains isotropy}
For a stable labelled type, the analytic moduli stack is
$[\mathcal T_{g,n}/P_{g,n}]$, with $\mathcal T_{g,n}$ contractible; the unordered stack is $[\mathcal T_{g,n}/M^+_{g,n}]$. Consequently its geometric \emph{stack} fundamental group is respectively $\widehat P_{g,n}$ or $\widehat{M^+_{g,n}}$. The manuscript's $M_{g,n}$ in (19)--(25), scans 326--329, corresponds to the unordered version: (20) explicitly has $\widehat{\mathcal T}^+_{g,n}$, not $\widehat{\mathcal T}^{!+}_{g,n}$. In particular
$\mathscr M_{0,3}=\Spec\mathbb Q$ but $\mathscr M_{0,[3]}\simeq B\mathfrak S_3$.
\Man{21 and 26}{128--131 and 165--166}{255--261 and 326--329}.
\begin{proof}[Verification]
The complex structure on an oriented surface gives the usual Teichm\"uller parameter; quotient by the identity component of the diffeomorphism group is the manuscript's $E_g/A^\circ_{g,n}$. Teichm\"uller contractibility, proved by Fenchel--Nielsen coordinates in \citet[Theorem 10.6 and \S10.6.3]{AF_FarbMargalit2012}, and the mapping class action give the stated quotient stack, retaining automorphism groups at its points. Analytic Riemann existence for these Deligne--Mumford stacks identifies finite covering spaces with finite \etale covers, hence gives their completed fundamental groups. One can prove this stack version by applying ordinary Riemann existence to an \etale atlas and its twofold and threefold intersections, then descending finite covers. This is the stack extension of the comparison used in \citet[\S13.9]{AF_BGH2020}. A labelled three-pointed projective line has one isomorphism class and no automorphisms; allowing permutations gives precisely $\mathfrak S_3$. A coarse-space quotient would erase this group, so it cannot satisfy the manuscript's sequence (20). The convention in STRATA F.16 is correct only with this unordered/labelled distinction retained explicitly.
\end{proof}

\statement{A.6}{Known}{Arithmetic fundamental-group and universal-monodromy sequences}
Fix $\overline{\mathbb Q}\hookrightarrow\mathbb C$ and compatible geometric base points on the unordered moduli stack and its universal punctured curve $\mathscr U$. For hyperbolic type,
\[
1\longrightarrow\widehat{M^+_{g,n}}\longrightarrow
 A^{\mathrm{arith}}_{g,n}:=\pi_1^{\et}(\mathscr M_{g,[n]})
 \longrightarrow G_{\mathbb Q}\longrightarrow1,
\]
\[
1\longrightarrow\pi_{g,n}\longrightarrow\pi_1^{\et}(\mathscr U)
 \longrightarrow A^{\mathrm{arith}}_{g,n}\longrightarrow1.
\]
Conjugation on the normal geometric fibre defines a commutative diagram to
\[
1\longrightarrow\pi_{g,n}\longrightarrow
 \Aut_{\mathrm{lac}}(\pi_{g,n})\longrightarrow O_{g,n}\longrightarrow1.
\]
The last sequence uses $Z(\pi_{g,n})=1$. The induced map
$\rho_{g,n}:A^{\mathrm{arith}}_{g,n}\to O_{g,n}$ is genuine after these choices; its change under base-point transport is conjugation in $O_{g,n}$. Its geometric restriction need not be identified with an injection before the congruence problem is settled.
\Man{26}{165--166}{326--329}.
\begin{proof}[Source and exactness]
The first sequence is the homotopy sequence of a geometrically connected stack over $\mathbb Q$. For the second sequence, the universal punctured curve is smooth with a smooth proper compactification and relative normal-crossings puncture divisor; its geometric fibres have the fixed hyperbolic type. The homotopy sequence is exact and the geometric fibre embeds; analytically it is the Birman sequence, and the profinite fibre injection is the assertion justified in Section C. Passing to an \etale atlas and descent gives the same sequence for the unordered stack. Normality makes conjugation a homomorphism into automorphisms, and centre-freeness identifies inner automorphisms with the fibre. Images of peripheral loops are determined by the puncture divisor, with the cyclotomic multiplier; this proves that the maps land in the indicated groups. Restricting to the algebraic closure recovers the completed topological monodromy. The labelled sequences are displayed explicitly in \citet[Corollary 6.4]{HoshiMochizuki2011}; the finite permutation quotient gives the unordered versions just proved. These are precisely (19)--(25) with all base points and orientation conventions stated.
\end{proof}

\subsection{A.3. Completion, discretification, and arithmetization}

\statement{A.7}{Proved}{A completion is not its image}
The discrete action induces a continuous surjection
\[
\widehat{M^+_{g,n}}\twoheadrightarrow
\Sigma_{g,n}:=\overline{\rho_{g,n}(M^+_{g,n})}\subset O_{g,n}.
\]
The kernel is the congruence kernel. The double hat in the manuscript denotes the ambient automorphism group $O_{g,n}$, not the profinite completion of $\widehat{M^+_{g,n}}$. In all formulae for stabilizers or normalizers inside $O_{g,n}$, it is the image $\Sigma$ that can be used without the congruence hypothesis.
\Man{26}{162r--163v}{320--323}.
\begin{proof}
The group $\pi_{g,n}$ is topologically finitely generated. Characteristic open subgroups are cofinal among its open normal subgroups, so $\Aut(\pi_{g,n})$ and $\Out(\pi_{g,n})$ have their usual profinite topologies. The discrete mapping class action is continuous for the topology pulled back from the finite outer automorphism groups, hence extends from its full profinite completion. Compactness makes its image closed, and density identifies that image with the closure of the discrete action. The universal property alone says nothing about the kernel. The manuscript explicitly writes on scan 320 that the natural surjection is being assumed injective. Section B supplies the exact known ranges and implications.
\end{proof}

\statement{A.8}{Proved}{The quotient sets behind the terminology}
Let $\pi$ be a profinite group with the loop structure of $\widehat\Pi_{g,n}$. Write
$I(\pi)=\Isom_{\mathrm{lac}}(\widehat\Pi_{g,n},\pi)$ and
$I^{\mathrm{out}}(\pi)=I(\pi)/\Inn(\pi)$.
Actual discretifications are the images of $\Pi_{g,n}$ under elements of $I(\pi)$, with their transported loop structures. Their based parameter set is
\[
\mathrm{Discr}(\pi)=I(\pi)/\Aut_{\mathrm{lac}}(\Pi_{g,n}).
\]
After taking inner-conjugacy classes, the denominator is $M^{\pm}_{g,n}$. Pre-discretifications replace this discrete denominator by its closure. If
$\Sigma^{\pm}=\overline{M^{\pm}_{g,n}}\subset O_{g,n}$, the outer quotient sets are
\[
\mathrm{Pre}(\pi)=I^{\mathrm{out}}(\pi)/\Sigma^{\pm},
\qquad
\mathrm{Pre}^{+}(\pi)=I^{\mathrm{out}}(\pi)/\Sigma.
\]
The maps from discretifications to pre-discretifications are quotient maps, not identifications of dense subgroups. Oriented pre-discretifications map two-to-one to unoriented ones for hyperbolic type.
\Man{26}{162--163}{320--323}; \S29, folios 188--189, scans 371--374.
\begin{proof}
Two identifications give the same embedded discrete group and transported structure exactly when their quotient restricts to an automorphism of the standard discrete group. An automorphism of the discrete group extends uniquely to its completion, which proves the first orbit description. Taking target inner conjugacy gives the outer description. The closure of the based discrete automorphism group contains $\Inn(\pi)$, because the discrete surface group is dense in $\pi$; thus the pre-discretification quotient is already outer. The sign character on the discrete mapping class group has image $\{\pm1\}$, and remains so on its closure. Its positive subgroup is closed of index two. This proves the last assertion without imposing the congruence property. It also explains the harmless change of notation from classes of discretifications in \S26 to pre-discretifications in \S29.
\end{proof}

\statement{A.9}{Proved}{Two different groups called $N$, and two different conjectures}
For a chosen geometric subgroup $\Sigma\subset O$, define
\[
N^{\mathrm{norm}}=N_O(\Sigma),\qquad
\mathfrak G^{\mathrm{norm}}=N^{\mathrm{norm}}/\Sigma.
\]
These are the \S26 definitions. In \S29 the manuscript defines instead
\[
N^{\mathrm{arith}}=\rho(A^{\mathrm{arith}}),\qquad
\mathfrak G^{\mathrm{arith}}=N^{\mathrm{arith}}/\Sigma.
\]
There is an inclusion $N^{\mathrm{arith}}\subseteq N^{\mathrm{norm}}$, and a canonical continuous surjection
$G_{\mathbb Q}\twoheadrightarrow\mathfrak G^{\mathrm{arith}}$ once the algebraic closure and the arithmetic marking have been fixed. The statement that this surjection is injective does not assert that $N^{\mathrm{arith}}=N^{\mathrm{norm}}$.
\Man{26 and 29}{164--165 and 188}{324--326 and 370--371}.
\begin{proof}
The geometric subgroup of the arithmetic fundamental group is normal. Taking its image under $\rho$ proves the inclusion in the normalizer. The image is compact and therefore closed. Quotienting the surjective map from $A^{\mathrm{arith}}$ to its image by the geometric subgroup gives the displayed surjection. No element outside the image is produced by this quotient operation. Hence replacing the arithmetic image by the full normalizer is an additional assertion, and changes an injectivity problem into a larger realization problem. Sections D and E settle the arithmetic injectivity and retain the other question separately.
\end{proof}

\noindent The dictionary now has an economical form:
\begin{center}\small
\begin{tabularx}{\textwidth}{@{}L{.27\textwidth}L{.27\textwidth}Y@{}}
\toprule Manuscript object & Precise modern object & Qualification\\\midrule
$\widehat{\mathcal T}^+$ & $\widehat{M^+}$ or its image $\Sigma$ & These agree exactly when the relevant congruence kernel vanishes.\\
$\widehat{\widehat{\mathcal T}}$ & $\Out_{\mathrm{lac}}(\pi)$ & Full continuous loop-preserving outer automorphisms.\\
$\operatorname{Discr\acute et}$ & Dense discrete surface models & Actual models and their inner-conjugacy classes are distinguished.\\
$\operatorname{Pr\acute ediscr\acute et}^+$ & $I^{\mathrm{out}}/\Sigma$ & A homogeneous space, not a group without additional normality.\\
Strict geometric structure & Conjugate of $\Sigma$ & Orbit parametrized by $O/N^{\mathrm{norm}}$.\\
Arithmetization $a$ & Class in $I^{\mathrm{out}}/N^{\mathrm{arith}}$ & Its oriented pre-discretifications form a $\mathfrak G^{\mathrm{arith}}$-torsor.\\
Arithmetic Teichm\"uller group & $\pi_1^{\et}(\mathscr M_{g,[n]})$ & Its image $N^{\mathrm{arith}}$ can be a congruence quotient.\\\bottomrule
\end{tabularx}
\end{center}

\statement{A.10}{Proved}{The arithmetization fibre is a torsor}
The map $I^{\mathrm{out}}/\Sigma\to I^{\mathrm{out}}/N^{\mathrm{arith}}$ has fibres that are torsors for $N^{\mathrm{arith}}/\Sigma$. This assertion needs normality of $\Sigma$ in the arithmetic image, already proved, but not normality in all of $O$.
\Man{29}{189--190}{373--376}.
\begin{proof}
For a representative $f$ of an arithmetization, its fibre consists of the right cosets $fn\Sigma$, for $n\in N^{\mathrm{arith}}$. Multiplication identifies it with a simply transitive $N^{\mathrm{arith}}/\Sigma$-set. Replacing $f$ changes this identification by translation, which is precisely why the fibre is canonically a torsor and has no canonical origin. After the faithfulness theorem in D, the acting group is $G_{\mathbb Q}$, subject to the fixed arithmetic identification.
\end{proof}

\subsection{A.4. Exceptional types and editorial limits}

\statement{A.11}{Proved}{The non-hyperbolic types do not satisfy the hyperbolic assertions}
The following table refers to punctures, not boundary circles fixed pointwise:
\begin{center}\small
\begin{tabular}{@{}llll@{}}\toprule
Type & $\Pi$ & $M^+$ & $P$\\\midrule
$(0,0)$ & $1$ & $1$ & $1$\\
$(0,1)$ & $1$ & $1$ & $1$\\
$(0,2)$ & $\mathbb Z$ & $\mathbb Z/2$ & $1$\\
$(0,3)$ & $F_2$ & $\mathfrak S_3$ & $1$\\
$(1,0)$ & $\mathbb Z^2$ & $\mathrm{SL}_2(\mathbb Z)$ & $\mathrm{SL}_2(\mathbb Z)$\\
$(1,1)$ & $F_2$ & $\mathrm{SL}_2(\mathbb Z)$ & $\mathrm{SL}_2(\mathbb Z)$\\\bottomrule
\end{tabular}
\end{center}
The extended groups are obtained by including degree $-1$ maps; for $(0,2)$ the result is $(\mathbb Z/2)^2$, and for both torus types it is $\mathrm{GL}_2(\mathbb Z)$. The surface group $F_2$ in the last line has peripheral element $[a,b]^{-1}$; the $F_2$ in type $(0,3)$ has three peripheral classes. They are not the same group with loops. Forgetting the marked point on a torus gives an isomorphism of mapping class groups, not a Birman sequence with an injected $\mathbb Z^2$ kernel.
\Man{19}{111--115}{222--230}; \S21, folios 129--131, scans 257--261.
\begin{proof}
A sphere with three ordered points has trivial orientation-preserving mapping class group, by cutting into discs; allowing their permutations gives $\mathfrak S_3$. Forgetting points gives the first three entries, with the remaining transposition for two unordered points. Complex conjugation fixes the real punctures and supplies the orientation-reversing involution; for two punctures it is independent of the permutation. A torus homeomorphism acts on its first homology by an integral matrix, and the linear maps realize every such matrix. Cutting along a basis proves that the action on homology detects the mapping class. A point can be fixed by translation, and translations are isotopic to the identity, so the pointed and unpointed groups coincide. For the cylinder, write $\rho_{\rm perm}$ for the permutation sign and $\rho_{\pi}$ for the sign on $\pi_1\simeq\Z$. Then $\rho_{\pi}=\rho_{\rm perm}\epsilon$. Indeed $z\mapsto z^{-1}$ has permutation sign $-1$ and orientation sign $+1$, while $z\mapsto\overline z$ has signs $+1,-1$; both induce $-1$ on $\pi_1$. Thus the identification of the two homomorphisms called $\rho$ on scan 225 is false if read literally, and the determinant-one image assertion for the rank-one case on scan 227 is false as well. The orientation-preserving cylinder group acts on $\Z$ by both signs. These computations also exhibit why the fundamental group of the identity component, which is nontrivial in the torus case, destroys the hyperbolic point-push injection. Their finite permutation and matrix checks are reproduced in \texttt{code/af\_checks.py}. For $(0,0)$ and $(0,1)$ the group alone has no orientation information; the manuscript adds an orientation line and the type, as stated explicitly at scans 229--230.
\end{proof}

\statement{A.12}{Proved}{Fields, real structures, and characteristics are part of the data}
The complex Teichm\"uller tower in this paper consists of orientable surfaces. A non-orientable surface is not the underlying surface of a smooth complex algebraic curve. For a real curve the relevant structure is an extension by $G_{\mathbb R}\simeq\mathbb Z/2$, together with its outer involution on the geometric group; it is not obtained by replacing the curve with the topological quotient by complex conjugation. All arithmetic comparison statements here are in characteristic zero. At primes $2$ and $3$, prime-to-characteristic or tame specialization cannot silently replace the full profinite group, and stack isotropy of orders $2$ and $3$ must be retained.
\begin{proof}
Holomorphic coordinate changes preserve the real orientation, proving the first assertion. The arithmetic fundamental-group sequence over $\mathbb R$ has geometric kernel and quotient $G_{\mathbb R}$; conjugation reverses the complex orientation and has cyclotomic value $-1$. Its quotient surface may have boundary or be non-orientable and is a different topological object. In characteristic $p$, wild inertia is pro-$p$ and is killed by prime-to-$p$ completion. Moreover elliptic curves in characteristic zero already have stabilizers of orders $2$, $4$, and $6$ at suitable points. Thus erasing all $2$- or $3$-primary isotropy cannot preserve the stack group. The $p$-adic statements cited in D have their own all-prime hypotheses; no comparison made here transports them to a curve over a finite field. These are scope restrictions, not an unproved assertion of a non-orientable analogue.
\end{proof}

\noindent\textbf{Reading record.}
On scan 204 the sequence of marked points is written with a positive-index convention, while (1) and the margin use indices beginning with zero. We use an arbitrary nested enumeration: the mathematics supports a reindexing, not a settled palaeographic correction. On scan 320 the assumption of injectivity is visible in the image and must not be suppressed. On scan 328 the transcription's folio alternatives $166|168$ do not alter any argument; the numerical sequence favours 166, without settling the reading. The phrase ``classes of discretifications'' at scan 321 is explicitly broadened by the insertion about pre-discretifications. At scan 261 the transcribed chain with changing genus cannot be literally correct by dimension: the intended zero-genus chain $\mathcal T_{0,3}\to\mathcal T_{0,2}\to\mathcal T_{0,1}\to\mathcal T_{0,0}$ fits the following sentence. That scan image was not supplied, so we record mathematical support only, not a visual emendation.
